\label{sec:discussion}

GLP-1 receptor agonists represent a major therapeutic advance for obesity, a chronic condition affecting approximately 40 percent of adult Medicaid enrollees. Despite proven effectiveness in treating obesity and type 2 diabetes, insurance coverage for obesity-indicated GLP-1 medications remains somewhat controversial. The drugs are expensive and most state Medicaid programs have opted not to cover GLP-1 medications for weight management. Several states that initially adopted coverage have since restricted or eliminated it, underscoring the fiscal pressures these drugs create. We exploited this policy variation to estimate the effects of coverage on utilization and spending in the Medicaid program.

Our findings indicate that Medicaid coverage substantially increases prescribing of obesity-indicated GLP-1 medications. In states that expanded formulary coverage, prescriptions increased by \rxObSevenZeroEmw{} per 100 enrollee-months in the first quarter of adoption. The effects grew over time: by nine quarters after adoption, prescriptions had increased by \rxObSevenNineEmw{} per 100 enrollee-months---equivalent to roughly 1 additional prescription per 120 enrollee-months, or per 10 enrollee-years. These rates are computed over \emph{all} Medicaid enrollees, not only those who might be clinically indicated for the drugs because of obesity or overweight. Given that approximately 40 percent of adult Medicaid enrollees meet clinical criteria for obesity and that the denominator includes children, the per-eligible-enrollee effect would be correspondingly larger. Coverage for obesity-indicated GLP-1 drugs did not meaningfully affect prescriptions for diabetes or cardiovascular indications, suggesting that coverage generated new utilization rather than substitution or relabeling of existing prescriptions. The results also indicate no positive spillovers in prescribing---no additional prescriptions for diabetes indications as a result of a newly covered (obesity) indication.

A related question is whether Medicaid coverage displaced private out-of-pocket spending on \emph{compounded GLP-1 products}, which grew rapidly over this period through online telehealth platforms that entered the market when key on-patent GLP-1s were in shortage. Our Consumer Edge analysis yields negative point estimates for compounding platform spending in states that adopted coverage, directionally consistent with some displacement, but the effects are small and not statistically significant. The utilization increases we document thus appear to reflect predominantly new access to GLP-1 obesity medications rather than substitution from lower priced compounded GLP-1 drugs.

A central challenge for assessing the budget implications of Medicaid coverage for obesity-indicated GLP-1 medications is that publicly available spending data reflect gross Medicaid reimbursements before manufacturer rebates. We addressed this constraint by estimating net spending using the statutory MDRP rebate formula and publicly available price proxies. Our estimates indicate that rebates reduce the spending effect by approximately \spReductionObSevenNineEmw{} percent relative to gross reimbursements. Although rebates are substantial, covering obesity-indicated GLP-1 medications represents a very large new outlay even on a net-of-rebate basis. 

These results are informative for ongoing debates about Medicaid coverage of GLP-1 obesity drugs. Several states that initially adopted coverage have since scaled back or eliminated it, citing budget pressures. Our estimates provide a basis for projecting the fiscal consequences of such decisions: if coverage effects are approximately symmetric, terminating coverage would reduce utilization and spending by magnitudes similar to those we estimate for adoption. Our findings are also relevant to proposals to expand Medicare Part D coverage of weight management medications. Although Medicaid and Medicare populations differ (Medicare enrollees are older and may have different obesity prevalence and comorbidity profiles, which could affect both uptake rates and the cost-effectiveness calculations), our estimates of coverage effects on utilization provide a starting point for projecting Medicare impacts. 

Our analysis has several limitations. First, we relied on a binary classification of coverage that does not distinguish variation in utilization management stringency. Most states that cover obesity-indicated GLP-1 medications have implemented prior authorization requirements with different criteria and enforcement stringency. Future research could examine how prior authorization stringency affects the utilization response to coverage \citep{liu2024state, klebanoff2025medicaid}. Second, we used aggregate state-quarter data from the SDUD rather than individual-level claims. This limits our ability to examine heterogeneity in coverage effects across enrollee subgroups defined by health status, demographics, or geography. It also precludes analysis of individual adherence patterns or health outcomes. Third, our net spending estimates rely on price proxies rather than actual rebate data. Validation against actual state fiscal data would strengthen confidence in the approach. Some states publish aggregate net-of-rebate pharmacy spending in budget documents, which could serve as an external benchmark, though differences in supplemental rebate inclusion and reporting conventions would complicate direct comparisons. Fourth, our analysis focuses on the direct fiscal costs of drug coverage and does not estimate potential offsets from reduced spending on obesity-related conditions. Recent evidence suggests that such offsets are limited in the short run \citep{wing2026glp1, zhang2026glp1, wennberg2025glp1}, but longer-run effects remain uncertain. Fifth, our projections extrapolate per-enrollee utilization at event time 9, which may not reflect steady-state prescribing if discontinuation rates are high or if new formulations alter long-run adherence patterns.

The policy landscape for GLP-1 obesity coverage continues to evolve.  Since July 2025, four states have eliminated Medicaid coverage of obesity-indicated GLP-1 medications: California\footnote{California's 2025-26 state budget eliminated Medi-Cal coverage of GLP-1s (Wegovy, Zepbound, Saxenda) for weight loss effective January 1, 2026. The state projected costs would have reached nearly \$800 million annually without the cut \citep{kffhealthnews2026california, medicalrx2025faq}.}, New Hampshire\footnote{New Hampshire ended Medicaid GLP-1 coverage for weight loss effective January 1, 2026 \citep{kff2026glp1medicaid}.}, Pennsylvania\footnote{Pennsylvania ended Medicaid coverage of GLP-1s for weight loss for adults aged 21 and older effective January 1, 2026. Coverage may continue for individuals under 21 under the EPSDT mandate. Pennsylvania spent \$650 million on GLP-1 prescriptions for Medicaid recipients in 2024, up from \$223 million in 2022 \citep{spotlightpa2025glp1, phlp2025glp1}.}, and South Carolina.\footnote{South Carolina eliminated Medicaid coverage of GLP-1s for weight loss effective January 1, 2026, reversing a policy that had been in place since November 2024 \citep{kff2026glp1medicaid}.} Michigan has substantially restricted its coverage, limiting eligibility exclusively to individuals classified as morbidly obese who have failed all other weight-loss interventions.\footnote{Michigan's FY2026 bipartisan state budget directed MDHHS to implement stricter criteria for GLP-1 medications prescribed for weight loss, effective January 1, 2026. Approval requires morbid obesity classification, documented failure of all other clinically appropriate weight-loss interventions and PDL-preferred anti-obesity agents, with coverage limited to averting the need for bariatric surgery \citep{mimeridian2025glp1, priorityhealth2025glp1}.} North Carolina temporarily eliminated coverage in October 2025 due to a legislative budget stalemate but reinstated it in December 2025 following a governor's directive.\footnote{North Carolina Medicaid ended GLP-1 coverage for obesity on October 1, 2025, citing shortfalls in state funding, but reinstated coverage on December 12, 2025, reverting to pre-October criteria \citep{ncmedicaid2025reinstate}.} Rhode Island, Wisconsin, and Connecticut are planning or considering restrictions.\footnote{See \citet{kffhealthnews2026california} and \citet{stateline2025retreat}. Connecticut's Governor Lamont proposed repealing the state's 2023 Medicaid GLP-1 coverage requirement amid a \$290 million Medicaid deficit \citep{ctmirror2025glp1}.} No states have newly added Medicaid coverage of GLP-1s for weight loss in 2026 thus far; KFF's 2025 Medicaid budget survey reported that state interest in expanding obesity drug coverage is waning, with cost cited as the primary barrier \citep{kff2026glp1medicaid}. Looking ahead, the federal BALANCE (Better Approaches to Lifestyle and Nutrition for Comprehensive hEalth) Model, announced by CMS in December 2025, will allow state Medicaid agencies to voluntarily opt in to cover GLP-1 medications for weight management at prices negotiated directly between CMS and participating manufacturers, launching as early as May 2026.\footnote{Under BALANCE, CMS negotiates drug pricing and coverage terms with GLP-1 manufacturers on behalf of participating state Medicaid agencies and Medicare Part D plan sponsors. Participation is voluntary for manufacturers, states, and plans \citep{cms2025balance_pr, cms2025balance_page}.} State decisions about coverage will depend on negotiated prices, evolving clinical evidence, and budget constraints. 